% One acme nut thread in axial section, at four stages of wear. Screw axis horizontal.
\begin{tikzpicture}[>=Latex,line join=round]
\tikzset{
  brz/.style={fill=orange!32,draw=orange!70!black,thick},
  body/.style={fill=orange!14,draw=orange!70!black,thick},
  gone/.style={fill=red!7,draw=red!50,thin,dashed},
  scr/.style={fill=blue!10,draw=blue!55!black,thick},
  ld/.style={->,red!70!black,line width=1.2pt},
  ttl/.style={font=\scriptsize\bfseries,align=center,text width=4.0cm,anchor=north},
  sub/.style={font=\tiny,align=center,text width=3.7cm,text=black!60,anchor=north}
}
\def\S{10.5}      % in -> cm, thickness scale
\def\Ht{1.25}     % tooth height (radial)
\def\tap{0.16}    % half-taper

% \tooth{x centre}{wear depth, in}{contact height fraction}{title}{subtitle}
\newcommand{\tooth}[5]{%
  \pgfmathsetmacro\tn{0.15*\S}            % new half-widths
  \pgfmathsetmacro\tw{(0.15-#2)*\S}       % worn thickness
  \pgfmathsetmacro\w{#2*\S}               % material lost, off the left flank
  % nut body above, bore to left and right
  \fill[body] (#1-2.10,\Ht) rectangle (#1+2.10,\Ht+0.42);
  % ghost of the original tooth
  \ifdim #2pt>0pt
    \fill[gone] (#1-\tn/2-\tap,0) -- (#1+\tn/2+\tap,0)
             -- (#1+\tn/2,\Ht) -- (#1-\tn/2,\Ht) -- cycle;
  \fi
  % the remaining tooth: left flank has retreated by \w
  \fill[brz] (#1-\tn/2-\tap+\w,0) -- (#1+\tn/2+\tap,0)
          -- (#1+\tn/2,\Ht) -- (#1-\tn/2+\w,\Ht) -- cycle;
  \draw[orange!70!black,thick] (#1-\tn/2-\tap+\w,0) -- (#1+\tn/2+\tap,0)
          -- (#1+\tn/2,\Ht) -- (#1-\tn/2+\w,\Ht) -- cycle;
  % the screw flank bearing on it, pushing left-to-right
  \fill[scr] (#1-2.10,-0.05) -- (#1-\tn/2-\tap+\w-0.06,-0.05)
          -- (#1-\tn/2+\w-0.06,\Ht) -- (#1-2.10,\Ht) -- cycle;
  \node[font=\tiny,text=blue!55!black,anchor=west] at (#1-2.05,\Ht*0.72) {screw};
  % load, applied over the part of the flank still in contact
  \pgfmathsetmacro\yc{\Ht*#3}
  \draw[ld] (#1-1.35,\yc) -- (#1-\tn/2-\tap+\w-0.10,\yc);
  % thickness at the tip (minor diameter), where the 0.15 in is quoted
  \draw[<->,gray!75!black,thin] (#1-\tn/2-\tap+\w,-0.42) -- (#1+\tn/2+\tap,-0.42);
  \node[font=\tiny,anchor=north,text=gray!50!black] at (#1,-0.44)
    {\pgfmathparse{0.15-#2}\pgfmathprintnumber[fixed,precision=3]{\pgfmathresult}\,in};
  \node[ttl] at (#1,-0.92) {#4};
  \node[sub] at (#1,-1.62) {#5};
}

\tooth{0}{0.000}{0.50}{New}{\SI{0.15}{in} at the minor diameter. Thread strength \textbf{over 15 times} the design load.}
\tooth{4.5}{0.037}{0.50}{At the replacement limit}{\SI{0.040}{in} end play. Still \textbf{10 times} stronger than the design load --- so the limit is \emph{not} a strength limit.}
\tooth{9.0}{0.080}{0.30}{Contact retreats}{``the tip of the thread stopped wearing while the rest of the face continued wearing''}
\tooth{13.5}{0.140}{0.16}{Sheared}{\SI{0.01}{in} left, about \SI{90}{\percent} gone. ``The fracture mode is shear.''}

\node[font=\footnotesize,align=center,text width=15.5cm,anchor=north] at (6.75,-2.95)
 {The bearing area falls with the thickness, so the pressure rises as $p\propto1/t$ --- and the wear
  rate rises with it. \textbf{The wear is not linear, and the last of it is the fastest.}};

\end{tikzpicture}
